
\documentclass[11pt]{article}
\usepackage[margin=1in]{geometry}
\usepackage{amsmath,amssymb,booktabs,hyperref,graphicx,parskip,enumitem}
\usepackage{xcolor}
\hypersetup{colorlinks=true,linkcolor=blue!50!black,urlcolor=blue!50!black}
\setlist{nosep}

\title{Replication Report: Common Mode Voltage Reduction of Single-Phase Quasi-Z-Source Inverter-Based Photovoltaic System\\[2pt]\large OSTI 1606674 \textbar{} Coverage 8/10, Agreement 9/10}
\author{Rick Stevens \and Ollie (AI Assistant)\\\small Argonne National Laboratory}
\date{2026-04-27}

\begin{document}\maketitle

\section{Source paper}
\begin{itemize}
\item \textbf{Title:} Common Mode Voltage Reduction of Single-Phase Quasi-Z-Source Inverter-Based Photovoltaic System
\item \textbf{Authors:} Sajadian, Ahmadi, Kouzani
\item \textbf{Year:} 2019
\item \textbf{Domain:} Power electronics
\item \textbf{Identifier:} OSTI 1606674
\end{itemize}


\section{Replication objective}
The central claim of the paper concerns power electronics. Our objective was to reproduce the headline numerical/qualitative results within the constraints of the AI-driven Replication Project (24--72\,h, commodity GPU/CPU compute, no author contact). A replication plan is archived at \texttt{1606674-Common-Mode-Voltage-Reduction-of-Single-Phase/replication\_plan.pdf}.

\section{Methodology}
We followed the standard Replication Project workflow: ingest the source PDF, extract method and data pointers, draft a replication plan, attempt the smallest end-to-end pipeline that produces a comparable number, then scale up where compute and data permit. Tools used in this replication are catalogued in the meta-paper's computational tools inventory (see \texttt{COMPUTATIONAL\_TOOLS\_INVENTORY.md}).



\subsection*{Paper-directory README excerpt}
\small
\par\medskip\textbf{Common Mode Voltage Reduction of Single-Phase Quasi-Z-Source Inverter-Based Photovoltaic System}\par
\par$\bullet$\ **OSTI ID:** 1606674
\par$\bullet$\ **Rank:** \#8
\par$\bullet$\ **Replication Score:** 9/10

\par\medskip\textbf{\# Why This Paper}\par
The paper presents a computational analysis of inverter modulation methods with clear equations, simulation outputs, and quantitative metrics, making it highly suitable for AI replication.

\par\medskip\textbf{\# Replication Plan}\par
AI would implement the inverter topology and modulation schemes in circuit simulation software, reproduce equations and switching states, and generate quantitative results for CMV and leakage current under various scenarios.

\par\medskip\textbf{\# Status}\par
\par$\bullet$\ [x] Paper reviewed
\par$\bullet$\ [x] Data identified
\par$\bullet$\ [x] Code implemented
\par$\bullet$\ [x] Results validated

\par\medskip\textbf{\# Replication Summary (2026-04-23, Ollie / OpenClaw)}\par

**Independent implementation:** `replication/simulate.py`
\par$\bullet$\ NumPy SPWM generator (unipolar, \$M=0.78\$, \$f\_s=10\$ kHz, shoot-through \$D=0.125\$) producing leg-voltage traces for both the traditional and the proposed (ZS1\$\textbackslash{}to\$ZS2) modulation.
\par$\bullet$\ PySpice / ngspice-46 transient analysis of the common-mode equivalent circuit (Eq.\textasciitilde{}5–6 of the paper) with \$L\_f=3\$ mH, \$C\_g=2.2\$ nF, \$R\_g=100\textbackslash{},\textbackslash{}Omega\$, \$R\_f=2\textbackslash{},\textbackslash{}Omega\$.

**Key numerical results:**
| Quantity | Traditional | Proposed | Reduction |
|---|---|---|---|
| \$\textbackslash{}lvert U\_\textbackslash{}text\{CM\}\textbackslash{}rvert\_\textbackslash{}text\{peak\}\$ | 400.0 V | 200.0 V | **50.0 \%** |
| \$U\_\textbackslash{}text\{CM,rms\}\$ | 223.1 V | 140.8 V | 36.9 \% |
| \$\textbackslash{}lvert I\_\textbackslash{}text\{LK\}\textbackslash{}rvert\_\textbackslash{}text\{peak\}\$ | 730 mA | 433 mA | 40.7 \% |
| \$I\_\textbackslash{}text\{LK,rms\}\$ | 269 mA | 170 mA | 36.7 \% |

The CMV-halving claim reproduces **exactly** (the proposed scheme clamps CMV to \$\textbackslash{}\{0,V\_\{PN\}/2\textbackslash{}\}\$ while the traditional scheme visits \$\textbackslash{}\{0,V\_\{PN\}/2,V\_\{PN\}\textbackslash{}\}\$). The switching-harmonic line at \$f\_s=10\$ kHz in the CMV spectrum drops by \$\textbackslash{}sim\$40 dB.

**Deliverables:**
\par$\bullet$\ `replication/simulate.py` — full simulation code
\par$\bullet$\ `replication/fig\_cmv\_ilk\_compare.png` — CMV + leakage current comparison (replicates Fig. 6(c) vs 8(c))
\par$\bullet$\ `replication/fig\_cmv\_spectrum.png` — CMV spectrum
\par$\bullet$\ `replication/metrics.txt` — quantitative metrics
\par$\bullet$\ `replication/report/replication\_report.pdf` — LaTeX report

**Score: 4 / 5** — central claim reproduced exactly; leakage-current ratio within the range consistent with unspecified CM-path parameters; no experimental quantities attempted.
\normalsize

The replication was driven by an OpenClaw agent loop (Claude Opus / GPT-5 class models) issuing shell, Python, and (where applicable) GPU jobs. Compute usage and wall-time for individual papers are summarised in the meta-paper's resource table; not separately recorded for this paper unless noted in the Tier-Lift artefact. Sub-agents were spawned for replication-plan generation and (for papers in batches 1--4) for tier-lift extension runs.

\section{What was reproduced}
\begin{itemize}
\item Unipolar SPWM at M=0.78, f\_s=10 kHz
\item Shoot-through duty D=0.125 qZSI modulation
\item Common-mode equivalent circuit (L\_f=3 mH, C\_g=2.2 nF, R\_g=100 Ohm, R\_f=2 Ohm)
\item Traditional vs proposed modulation leg-voltage traces
\item 50\% CMV peak reduction (exact)
\item \textasciitilde{}40 dB CMV switching-harmonic attenuation at 10 kHz
\item Leakage-current waveform comparison (Fig. 6c vs 8c)
\end{itemize}

\section{What was missing or incomplete}
\begin{itemize}
\item Experimental hardware prototype
\item MPPT and grid-synchronisation control loops
\item Full EMI filter stage
\item Efficiency and THD numbers
\item Component-tolerance sensitivity study
\item Thermal / conduction-loss analysis
\end{itemize}
The dominant blocker categories for this paper, mapped onto the meta-paper's 9-category friction taxonomy (\S4), are listed in \S\ref{sec:friction} below.

\section{Quantitative agreement}
Per-quantity numerical comparisons are not separately tabulated in the structured evaluation record for this paper beyond the agreement rationale below; where the rationale references specific numbers, they are reproduced verbatim. A full numerical comparison table is recorded in the per-replication artefacts (see \S\ref{sec:repro}). For papers below 8/8, the agreement rationale identifies the specific quantities that diverge.

\paragraph{Agreement rationale.} CMV halving reproduces exactly: |U\_CM|\_peak 400 V -\textgreater{} 200 V (50\% reduction), matching the paper's central claim that proposed scheme clamps CMV to \{0, V\_PN/2\} while traditional scheme visits \{0, V\_PN/2, V\_PN\}. Leakage current peak 40.7\% reduction (730 mA -\textgreater{} 433 mA) and RMS 36.7\% reduction align with the paper's reported envelope (paper reports leakage in the same ballpark; exact match not expected without full CM-path parasitics). Switching-harmonic line at 10 kHz drops by \textasciitilde{}40 dB in our spectrum, consistent with paper Fig. 8(c).

\section{Score and friction-point tagging}\label{sec:friction}
\begin{itemize}
\item \textbf{Coverage:} 8/10 --- Independent NumPy SPWM generator plus PySpice/ngspice transient simulation of the common-mode equivalent circuit (Eqs 5-6 of the paper) with the same modulation indices, switching frequency, shoot-through duty, and CM-path R/L/C values. Proposed ZS1-\textgreater{}ZS2 scheme and traditional scheme both implemented. Not replicated: hardware prototype measurements, MPPT loop, full dq grid-synchronisation block, EMI filter design, component-tolerance Monte Carlo.
\item \textbf{Agreement:} 9/10 --- CMV halving reproduces exactly: |U\_CM|\_peak 400 V -\textgreater{} 200 V (50\% reduction), matching the paper's central claim that proposed scheme clamps CMV to \{0, V\_PN/2\} while traditional scheme visits \{0, V\_PN/2, V\_PN\}. Leakage current peak 40.7\% reduction (730 mA -\textgreater{} 433 mA) and RMS 36.7\% reduction align with the paper's reported envelope (paper reports leakage in the same ballpark; exact match not expected without full CM-path parasitics). Switching-harmonic line at 10 kHz drops by \textasciitilde{}40 dB in our spectrum, consistent with paper Fig. 8(c).
\end{itemize}

\noindent\textbf{Friction points encountered (taxonomy tags):}
\begin{itemize}
\item None recorded.
\end{itemize}

\section{Follow-on questions}
\begin{itemize}
\item Q1: How does the 50\% CMV reduction degrade as the CM parasitic capacitance C\_g varies over 0.5-10 nF (representative of different PV-array sizes and mounting conditions)?
\item Q2: Does the ZS1-\textgreater{}ZS2 modulation preserve its CMV advantage under dead-time insertion of 0.5-2 us, or does dead-time reintroduce intermediate CM levels?
\item Q3: Can a third modulation state (ST-\textgreater{}ZS3) further reduce CMV peak below V\_PN/4 without penalising DC-link utilisation, and what is the achievable theoretical floor?
\item Q4: How do the CMV and leakage-current reductions translate to measured conducted EMI spectra (150 kHz-30 MHz) on a prototype, and does a smaller EMI filter become viable?
\item Q5: When the grid voltage is distorted (\textgreater{}5\% THD, notched by nonlinear loads), does the proposed scheme maintain the clean two-level CMV, or does the control interaction reintroduce transients?
\end{itemize}

\section{Reproducibility pointers}\label{sec:repro}
\begin{itemize}
\item \textbf{Paper directory:} \texttt{1606674-Common-Mode-Voltage-Reduction-of-Single-Phase}
\item \textbf{Replication artefacts:}
\begin{itemize}
\item \texttt{/Users/stevens/Dropbox/REPLICATE-PROJECT/1606674-Common-Mode-Voltage-Reduction-of-Single-Phase/replication}
\end{itemize}
\item \textbf{Replication-dir hint from JSONL:} \texttt{\textasciitilde{}/Dropbox/REPLICATE-PROJECT/1606674-Common-Mode-Voltage-Reduction-of-Single-Phase/replication/}
\item \textbf{Status:} Not recorded
\end{itemize}

To re-run, change directory into the paper directory above and consult its \texttt{README.md} for the entry-point script. Most replications follow the convention \texttt{replication/run\_*.sh} or \texttt{replication/<subdir>/run.py}.

\end{document}
